\begingroup
\renewcommand{\arraystretch}{1.45}
\fancyhead[L]{\small ARITHMETIC–GEOMETRIC MEAN}
\Needspace{43\baselineskip}
\AGMArticle{8}
If we set
\[
1+\tfrac14xx+\tfrac14.\tfrac9{16}x^4
+\tfrac14.\tfrac9{16}.\tfrac{25}{36}x^6+\text{etc.}=y
\]
then
\[
\begin{aligned}&\tfrac12xx+\tfrac14.\tfrac94x^4\\&\quad+\tfrac14.\tfrac9{16}.\tfrac{25}{6}x^6\\&\quad+\tfrac14.\tfrac9{16}.\tfrac{25}{36}.\tfrac{49}{8}x^8
+\text{etc.}=\frac{x\mathrm{d}y}{\mathrm{d}x}\end{aligned}
\]
and
\[
\begin{aligned}&xx+\tfrac14.9x^4+\tfrac14.\tfrac9{16}.25x^6\\&\quad+\tfrac14.\tfrac9{16}.\tfrac{25}{36}.49x^8+\text{etc.}\\&=\frac{xx\mathrm{dd}y}{\mathrm{d}x^2}+\frac{x\mathrm{d}y}{\mathrm{d}x}\end{aligned}
\]
whence it follows directly that
\[
\frac{xx\mathrm{dd}y}{\mathrm{d}x^2}+3\frac{x\mathrm{d}y}{\mathrm{d}x}+y
=\frac1{xx}\left(xx\frac{\mathrm{dd}y}{\mathrm{d}x^2}+\frac{x\mathrm{d}y}{\mathrm{d}x}\right)
\]
or
\[
(x^3-x)\frac{\mathrm{dd}y}{\mathrm{d}x^2}
+(3xx-1)\frac{\mathrm{d}y}{\mathrm{d}x}+xy=0
\]
In this way, therefore, our arithmetic–geometric means have been reduced to integral quantities
and supply a particular solution of this differential equation of the second order.

\Needspace{12\baselineskip}
The complete integral of the same equation is
\(\dfrac{\mathfrak{A}}{\mathrm{M}(1+x,\,1-x)}+\dfrac{\mathfrak{B}}{\mathrm{M}(1,x)}\).

\Needspace{32\baselineskip}
Let \(\varphi\) be an arbitrary angle. The value of the integral
\(\int\cos\varphi^2\mathrm{d}\varphi\), from \(\varphi=0\) to
\(\varphi=\pi\), will be, as is commonly known, \(=\tfrac12\pi\); in the same way,
the value of the integral \(\int\cos\varphi^4\mathrm{d}\varphi\) between the same limits is
\(=\tfrac12.\tfrac34\pi\); the value of the integral is given by
\(\int\cos\varphi^6\mathrm{d}\varphi=\tfrac12.\tfrac34.\tfrac56\pi\)
etc. --- finally, as is self-evident, \(\int\mathrm{d}\varphi=\pi\).
Hence it is clear that the value of the integral
\[
\begin{aligned}&\int\mathrm{d}\varphi\times
\bigl(1+\tfrac12x^2\cos\varphi^2\\&\quad+\tfrac12.\tfrac34x^4\cos\varphi^4\\&\quad+\tfrac12.\tfrac34.\tfrac56x^6\cos\varphi^6+\text{etc.}\bigr)\end{aligned}
\]
or of this one, \(\int\dfrac{\mathrm{d}\varphi}{\surd(1-xx\cos\varphi^2)}\),
is \(=\pi y\) if it is taken from \(\varphi=0\) to \(\varphi=\pi\),
regarding the quantity \(x\) as constant.

\Needspace{25\baselineskip}
Now if the function \(\dfrac1{\surd(1-xx\cos\varphi^2)}\) is supposed to be expanded into a series of the form
\[
P+2Q\cos2\varphi+2R\cos4\varphi+2S\cos6\varphi+\text{etc.}
\]
so that the coefficients \(P,Q,R,S\) etc. depend on \(x\) alone, the complete value
of the integral given above will be
\[
P\varphi+Q\sin2\varphi+\tfrac12R\sin4\varphi+\tfrac13S\sin6\varphi
+\text{etc.}+\text{Const.}
\]
\SourceLeaf{380}
\Needspace{31\baselineskip}
and thus the value within the limits previously stated will be $=P\varphi$, whence $y=P$. We now observe
that, since $\frac{1}{y}=\mathrm{M}(1+x,1-x)$, it is also the case that
$\frac{1}{y}=\mathrm{M}(1,\surd(1-xx))$. From this the following more general theorem is very readily derived:
If an expression of the form
$\frac{\alpha}{\surd(\beta-\gamma\cos\varphi^{2})}=W$
is expanded into a series proceeding according to cosines of the angles $2\varphi$, $4\varphi$, $6\varphi$ etc.,
and its constant term is denoted by $P$, while the maximum
and minimum values of $W$ (namely $\frac{\alpha}{\surd(\beta-\gamma)}$ and
$\frac{\alpha}{\surd\beta}$) are denoted by $v,v'$, then $\frac{1}{P}$ will be the arithmetic–geometric mean between $\frac{1}{v}$ and $\frac{1}{v'}$. In the first instance,
the preceding arguments apply only to the case where $\gamma$ is a positive
quantity; but they extend very readily also to the case where $\gamma$ is negative.
For in this case the result is $W=\frac{\alpha}{\surd(\beta-\gamma+\gamma\cos\psi^{2})}$
on setting $\psi=90-\varphi$; whence $\frac{1}{P}$ is the mean between
$\frac{\surd\beta}{\alpha}$ and $\frac{\surd(\beta-\gamma)}{\alpha}$. As above.
\Needspace{28\baselineskip}
Moreover, it is readily seen that the same theorem extends without any change also to expressions of the form $\frac{\alpha}{\surd(\beta+\gamma\cos\varphi)}$, namely, that
$\frac{1}{\text{Const. term}}$ becomes
\[
\begin{gathered}\\=\mathrm{M}\left(\frac{1}{\text{max. value}},\frac{1}{\text{min. value}}\right)\\=\mathrm{M}\left(\frac{\surd(\beta-\gamma)}{\alpha},
                 \frac{\surd(\beta+\gamma)}{\alpha}\right)\end{gathered}
\]
for functions of this kind will be reduced to the form
$\frac{\alpha}{\surd(\beta-\gamma+2\gamma\cos\psi^{2})}$, entirely similar to the one of which we have just spoken,
if we write $\varphi=2\psi$\HistoricalNote{29}.

\Needspace{24\baselineskip}
Finally, we point out that a much more general proof of these same theorems
will be given below from more natural principles; moreover, we shall soon also show how to obtain all
the remaining coefficients $Q,R,S$ etc. by equally convenient methods.
Here this inquiry will also serve to make these investigations dearer to those who,
while having no taste for the sublimity and divine beauty of eternal truths, know
how to judge their worth solely by the use they can bring to the branches of
applied mathematics. Everyone knows how useful
an expansion of the coefficients $P,Q,R,S$ etc. as rapid as the one that follows from these principles
is in physical astronomy, or the theory of planetary perturbations.


\par\begin{center}\rule{.15\linewidth}{.4pt}\end{center}
\clearpage
\endgroup
